\section{Ôn tập chương 2: Dãy số - cấp số}
\Opensolutionfile{ans}[ans/ans-0-B15]
\caulc
 \begin{ex}%Câu 1%[1D2N1-2]
 	Cho dãy số $\left(u_n\right)$ với $u_n=11^n.$ Tìm số hạng $u_{n+1}$.
 	\choice
 	{\True $u_{n+1}=11\cdot11^n$}
 	{$u_{n+1}=11^n+11$}
 	{$u_{n+1}=11\left(n+1\right)$}
 	{$u_{n+1}=11^n+1$}
 	\loigiai{
 Ta có 
 		$u_{n+1}=11^{n+1}=11\cdot11^n.$}
 \end{ex}


\begin{ex}%Câu 2%[1D2H1-4]
	Cho dãy số $\left(u_n\right)$ với $u_n=-7n+4$. Khẳng định nào dưới đây \textbf{sai}?
	\choice
	{Dãy số bị chặn trên}
	{Dãy số không bị chặn}
	{Dãy số giảm}
	{\True Dãy số tăng}
	\loigiai{
		\begin{itemize}
			\item Ta có $u_n=-7n+4\Rightarrow{u_{n+1}}-u_n=\left[-7(n+1)+4\right]-\left[-7n+4\right]=-7 < 0,\forall n.$ \\
			Suy ra, $\left(u_n\right)\,$ là dãy số giảm.
		\item Mà $u_n=-7n+4\le -7\cdot1+4=-3, \forall n\in\mathbb{N}^*$ nên dãy số bị chặn trên.
		\item Mặt khác, khi $n$ càng lớn thì $u_n$ càng nhỏ, nên dãy số không bị chặn dưới, do đó dãy số không bị chặn.
\end{itemize}	
}
\end{ex}


\begin{ex}%Câu 3%[1D2H1-5]
	Xét tính bị chặn của dãy số $u_n=-n^2+6n-7$.
	\choice
	{Bị chặn}
	{Không bị chặn}
	{\True Bị chặn trên}
	{Bị chặn dưới}
	\loigiai{
		Ta có $u_n=-n^2+6n-7=2-\left(n+3\right)^2< 7$.
		\begin{itemize}
			\item  Dãy số $\left(u_n\right)$ bị chặn trên; 
			\item Dãy dãy $\left(u_n\right)$ không bị chặn dưới.
	\end{itemize}	
	}
\end{ex}

\begin{ex}%Câu 4%[1D2N1-2]
	Cho dãy số $\left(u_n\right)$ có các số hạng đầu là $\dfrac{1}{5}$, $\dfrac{1}{5^2}$, $\dfrac{1}{5^3}$, $\dfrac{1}{5^4}$, $\dfrac{1}{5^5}$, $\cdots $ Số hạng tổng quát của dãy số này là
	\choice
	{$u_n=\dfrac{1}{5}.\dfrac{1}{5^{n+1}}$}
	{$u_n=\dfrac{1}{5^{n+1}}$}
	{\True $u_n=\dfrac{1}{5^n}$}
	{$u_n=\dfrac{1}{5^{n-1}}$}
	\loigiai{
		Dãy $u_n=\dfrac{1}{5^n}$ có khai khiển là $\dfrac{1}{5}$, $\dfrac{1}{5^2}$, $\dfrac{1}{5^3}$, $\dfrac{1}{5^4}$, $\dfrac{1}{5^5}$, $\cdots$
		}
\end{ex}

\begin{ex}%Câu 5%[1D2N2-2]
	Cho cấp số cộng $\left(u_n\right)$ có số hạng tổng quát là $u_n=5n-1$. Tìm công sai $d$ của cấp số cộng.
	\choice
	{\True $d=5$}
	{$d=6$}
	{$d=-1$}
	{$d=4$}
	\loigiai{
Ta có 	$ d=u_n-u_{n-1}=5$.
}
\end{ex}


\begin{ex}%Câu 6%[1D2H2-2]
	Cho cấp số cộng $\left(u_n\right)$ có $u_1=2024$ và công sai $ d=-20$. Tìm số hạng $u_{100}$.
	\choice
	{$u_{100}=2024\cdot \left(-20\right)^{99}$}
	{\True $u_{100}=44$}
	{$u_{100}=924$}
	{$u_{100}=946$}
	\loigiai{
	Ta có	$u_{100}=u_1+99\cdot d=2024+99\cdot\left(-20\right)=44$.}
\end{ex}

\begin{ex}%Câu 7%[1D2H2-5]
Cho cấp số cộng $\left(u_n\right)\colon 3$, $a$, $13$, $b$. Tích $ab$ bằng
	\choice
	{\True $144$}
	{$39$}
	{$26$}
	{$104$}
	\loigiai{
		Ta có $\heva{&	3+13=2a\\&	a+b=2.13
	}\Leftrightarrow\heva{&	a=8\\&b=18
	}\Rightarrow ab=144$.
}
\end{ex}


\begin{ex}%Câu 8%[1D2H2-7]
	Trong dịp tết trồng cây, một trường THPT muốn trồng $50$ hàng cây theo quy cách như sau:
	\begin{itemize}
		\item Hàng 1: Trồng $2$ cây.
	\item Hàng 2: Trồng $4$ cây.
	\item Hàng 3: Trồng $6$ cây.\\
	$\cdots \cdots$	
\end{itemize}
	Hỏi nhà trường trồng được bao nhiêu cây?
	\choice
	{\True $ 2550$}
	{$ 2650$}
	{$ 5300$}
	{$ 2500$}
	\loigiai{
		Ta có $u_1=2$, $d=2$ nên $S_{50}=50\cdot2+\dfrac{50\cdot \left(50-1\right).2}{2}=2550$.}
\end{ex}


\begin{ex}%Câu 9%[1D2N3-2]
	Trong các dãy số cho dưới đây, dãy số nào là cấp số nhân?
	\choice
	{$1$; $2$; $3$; $4$; $5$}
	{$0$; $4$; $8$; $12$; $16$}
	{$1$; $3$; $5$; $7$; $9$}
	{\True $ 5$; $\dfrac{5}{2}$; $\dfrac{5}{4}$; $\dfrac{5}{8}$; $\dfrac{5}{16}$}
	\loigiai{
Ta thấy ở dãy số $ 5$; $\dfrac{5}{2}$; $\dfrac{5}{4}$; $\dfrac{5}{8}$; $\dfrac{5}{16}$ có $u_1=5$, $u_2=\dfrac{5}{2}$, $u_3=\dfrac{5}{4}$, $u_4=\dfrac{5}{8}$, $u_5=\dfrac{5}{16}$ nên đây là cấp số nhân với công bội $ q=\dfrac{1}{2}$.}
\end{ex}

\begin{ex}%Câu 10%[1D2N3-2]
	Cho cấp số nhân $\left(u_n\right)$ với $u_1=3$ và $u_2=15$. Công bội của cấp số nhân đã cho bằng
	\choice
	{\True $ 5$}
	{$ 12$}
	{$\dfrac{1}{5}$}
	{$ 4$}
	\loigiai{
		Ta có $u_2=u_1\cdot q\Rightarrow q=\dfrac{u_2}{u_1}=\dfrac{15}{3}=5$.}
\end{ex}

\begin{ex}%Câu 11%[1D2H3-4]
	Cho cấp số nhân có $u_1=-8$, $q=\dfrac{3}{2}$. Số $-\dfrac{2187}{16}$ là số hạng thứ mấy của cấp số này?
	\choice
	{\True Thứ $8$}
	{Không phải là số hạng của cấp số}
	{Thứ $9$}
	{Thứ $7$}
	\loigiai{
Giả sử $u_n=-\dfrac{2187}{32}$.\\
		 Khi đó $$u_n=u_1q^{n-1}\Rightarrow-\dfrac{2187}{16}=\left(-8\right)\cdot{\left(\dfrac{3}{2}\right)^{n-1}}\Rightarrow{\left(\dfrac{3}{2}\right)^{n-1}}=\dfrac{2187}{128}=\left(\dfrac{3}{2}\right)^7.$$
		Suy ra $n-1=7\Rightarrow n=8$.\\
Vậy $-\dfrac{2187}{16}$ là số hạng thứ $8$.
}
\end{ex}

\begin{ex}%Câu 12%[1D2H3-5]
Xác định $ x$ để $3$ số $2x-3$; $x$; $2x+3$ theo thứ tự lập thành một cấp số nhân.
	\choice
	{$x=\pm\dfrac{1}{3}$}
	{\True $x=\pm\sqrt3$}
	{$x=\pm\dfrac{1}{\sqrt 3}$}
	{$x=\pm 3$}
	\loigiai{
Theo tính chất cấp số nhân $\left(2x-3\right)\left(2x+3\right)=x^2\Leftrightarrow 4x^2-9=x^2\Leftrightarrow x=\pm\sqrt 3 $.
}
\end{ex}
\Closesolutionfile{ans}
\indapan{6}{ans/ans-0-B15}

\cauds
\Opensolutionfile{ans}[ans/ans-0-B15-DS]
\begin{ex}%Câu 13%[1D2H2-6]
	Cho cấp số cộng $\left(u_n\right)$có $u_n=5n-8$. Xác định tính đúng, sai của các phát biểu sau đây
	\choiceTF
	{\True $u_1=-3$}
	{Công sai $ d=-8$}
	{\True Số $492$ là số hạng thứ $100$ của $\left(u_n\right)$}
	{Tổng của $10$ số hạng đầu tiên của cấp số cộng bằng $ 190$}
\loigiai{
\begin{itemchoice}
	\itemch $u_1=5\cdot1-8=-3$
	\itemch $ d=u_2-u_1=2-\left(-3\right)=5$.
	\itemch $u_{100}=5\cdot100-8=492$.
	\itemch $S_{10}=\dfrac{10}{2}\left(2u_1+9d\right)=5\left[2\left(-3\right)+9\cdot5\right]=195$.
\end{itemchoice}
}
\end{ex}

\begin{ex}%Câu 14%[1D2H3-6]
Cho cấp số nhân $\left(u_n\right)$ có $u_1=2$; $u_2=-4$. Xác định tính đúng, sai của các phát biểu sau đây
\choiceTF
{Công bội $ q=2$}
{$u_5=-32$}
{\True Số $-64$ là số hạng thứ $6$ của $\left(u_n\right)$}
{\True Tổng của $8$ số hạng đầu tiên của cấp số nhân bằng $-170$}
\loigiai{
\begin{itemchoice}
\itemch $ q=\dfrac{u_2}{u_1}=\dfrac{-4}{2}=-2$.
\itemch $u_5=u_1\cdot q^4=2\cdot\left(-2\right)^4=32$.
\itemch $u_6=u_1\cdot q^5=2\cdot\left(-2\right)^5=-64$.
\itemch $S_8=\dfrac{u_1\left(1-q^8\right)}{1-q}=\dfrac{2\left[1-\left(-2\right)^8\right]}{1-\left(-2\right)}=-170$.
\end{itemchoice}
}
\end{ex}

\begin{ex}%Câu 15%[1D2H2-6]
Cho cấp số cộng $\left(u_n\right)$ có số hạng đầu $u_1=-5$ và công sai $ d=3$.
\choiceTF
{\True Số $100$ là số hạng thứ $36$ của cấp số cộng}
{ Số hạng thứ $5$ của cấp số cộng là $ 9$}
{\True Tổng $10$ số hạng đầu của cấp số cộng là $85$}
{Số hạng tổng quát cảu cấp số cộng $\left(u_n\right)$ là $\left(u_n\right)=3n-7$}
\loigiai{
\begin{itemchoice} 
\itemch Ta có
$u_n=u_1+(n-1)d\Leftrightarrow 100=-5+(n-1)3\Leftrightarrow n=36$.
\itemch Ta có $u_5=u_1+4d=-5+4\cdot3=7$.
\itemch Ta có
$S_n=\dfrac{n}{2}\left[2u_1+(n-1)d\right]\Rightarrow{S_{10}}=\dfrac{10}{2}\left[2\cdot(-5)+(10-1)3\right]=85$.
\itemch $u_n=u_1+\left(n-1\right)d=-5+\left(n-1\right)3=3n-8$.
\end{itemchoice}}
\end{ex}

\begin{ex}%Câu 16%[1D2H3-6]
Cho cấp số nhân $\left(u_n\right)$ có tổng $n$ số hạng đầu tiên là $S_n=5^n-1$ với $n=1,2,\ldots$.
\choiceTF
{Công bội của cấp số nhân đã cho là $q=6$ }
{ \True Số hạng đầu của cấp số nhân đã cho là $u_1=4$}
{\True Tổng $5$ số hạng đầu của cấp số nhân là $ 3124$}
{Số hạng thứ $5$ của cấp số nhân là $ 2600$}
\loigiai{
\begin{itemchoice}
\itemch Ta có $\heva{&{u_1}=S_1=5-1=4\\&{u_1}+u_2=S_2=5^2-1=24
}\Rightarrow \heva{&{u_1}=4\\&{u_2}=24-u_1=20
.}$\\
Do đó  $u_1=4$, và $q=\dfrac{u_2}{u_1}=5$.
\itemch $u_1=4$.
\itemch $S_5=5^5-1=3124$.
\itemch Ta có $u_5=u_1.q^4=4\cdot5^4=2500$.
\end{itemchoice}
}
\end{ex}

\Closesolutionfile{ans}
\indapan{3}{ans/ans-0-B15-DS}

\Opensolutionfile{ans}[ans/ans-0-B15-KQ]
\caukq
 \begin{ex}%Câu 17%[1D2V2-4]
 Một cấp số cộng có $7$ số hạng. Biết rằng tổng của số hạng đầu và số hạng cuối bằng $30$, còn tổng của số hạng thứ ba và số hạng thứ sáu bằng $35$. Tính số hạng thứ bảy của cấp số cộng đó.
 \shortans{$30$}
 \loigiai{
Ta có: $\heva{&u_1+u_7=30\\&u_3+u_6=35}
 		 \Leftrightarrow \heva{&2u_1+6d=30\\&2u_1+7d=35}
 \Leftrightarrow \heva{&u_1=0\\&d=5.}$\\
 	Vậy $u_7=u_1+6d=0+6\cdot5=30$.
 }
\end{ex}

\begin{ex}%Câu 18%[1D2V2-7]
Người ta trồng cây theo hình tam giác, với quy luật: ở hàng thứ nhất có $1$ cây, ở hàng thứ hai có $2$ cây, ở hàng thứ ba có $3$ cây,…ở hàng thứ $n$ có $ n$ cây. Biết rằng người ta trồng hết $4950$ cây. Hỏi số hàng cây được trồng theo cách trên là bao nhiêu?
\shortans{$99$}
\loigiai{
Cách trồng cây theo quy luật trên lập thành một cấp số cộng với số hạng đầu là số cây ở hàng 1, $u_1=1$ và công sai $ d=1$.\\
Ta có 
\allowdisplaybreaks
\begin{eqnarray*}
&&S_n=4950 \Leftrightarrow 4950=\dfrac{n}{2}\left[2u_1+(n-1)d\right]\\
&\Leftrightarrow& 4950=\dfrac{n}{2}\left[2\cdot1+(n-1)\cdot 1\right]\\
	&\Leftrightarrow&{n^2}+n-9900=0\Leftrightarrow \hoac{&n=99\\&n=-100\,(\text{loại}).}
 \end{eqnarray*}
Vậy để trồng hết $4950$ cây thì cần $99$ hàng cây.

}
\end{ex}

\begin{ex}%Câu 19%[1D2V3-8]
Ông A gửi $ 120$ triệu đồng vào ngân hàng với lãi suất $ 6\% /$năm. Biết rằng nếu không rút tiền ra khỏi ngân hàng thì cứ sau mỗi năm số tiền lãi sẽ được nhập vào vốn để tính lãi cho năm tiếp theo. Hỏi sau $ 10$ năm, tổng số tiền mà ông $ A$ nhận được là bao nhiêu, giả định trong khoảng thời gian này lãi suất không thay đổi và ông $ A$ không rút tiền ra (kết quả làm tròn đến hàng đơn vị: triệu đồng)?
\shortans{$215$}
\loigiai{
Ta có số tiền ban đầu ông A có được là $u_1=120$.\\
Số tiền sau 1 năm ông A nhận được là $u_2=u_1+6\%\cdot u_1=u_1\cdot 1{,}06$.\\
Số tiền sau 2 năm ông A nhận được là $u_3=u_2+6\% \cdot u_2=u_2\cdot 1{,}06$.\\
Tương tự cho các năm tiếp theo.\\
Do đó, ta có số tiền gốc và lãi hàng năm ông A nhận được lập thành một cấp số nhân với $u_1=120$ triệu đồng và công bội $q=1{,}06$. \\
Tổng số tiền gốc và lãi sau 10 năm chính là số hạng thứ $11$ của cấp số nhân này và bằng $u_{11}=u_1\cdot q^{10}=120\cdot\left(1{,}06\right)^{10}\approx 215$ triệu đồng.
}
\end{ex}

\begin{ex}%Câu 20%[1D2V3-6]
Cho cấp số nhân$\left(u_n\right)$ có $ \heva{&u_3=16\\&u_2+u_4=40}$. Tính tổng $8$ số hạng đầu của cấp số nhân $\left(u_n\right)$, biết $ q > 1$.
\shortans{$1020$}
\loigiai{
Ta có $\heva{&u_3=16\\&u_2+u_4=40}\Leftrightarrow\heva{&u_1\cdot q^2=16\\&u_1\cdot q+u_1\cdot q^3=40.}$\\
Suy ra $\dfrac{q}{1+q^2}=\dfrac{16}{40}\Leftrightarrow 2q^2-5q+2=0\Leftrightarrow\hoac{&q=2\\&q=\dfrac{1}{2}.}$\\
Vì $ q > 1$ nên $ q=2$. Khi đó $u_1=4$.\\
Tổng 8 số hạng đầu của cấp số nhân $\left(u_n\right)$ là $S_8=\dfrac{u_1\left(1-q^8\right)}{1-q}=\dfrac{4\left(1-2^8\right)}{1-2}=1020$.
}
\end{ex}

\begin{ex}%Câu 21%[1D2V3-4]
Cho hình vuông $A_1$ có cạnh bằng $ 1$. Gọi $A_2$ là hình vuông có các đỉnh là trung điểm các cạnh của hình vuông $A_1$; $A_3$ là hình vuông có các đỉnh là trung điểm các cạnh của hình vuông $A_2$,... Cứ tiếp tục quá trình như trên ta được dãy các hình vuông $A_1;\,A_2;\,...;\,A_n;...$ Biết rằng diện tích của hình vuông $A_{2024}$ là $S=\dfrac{1}{2^{\overline{abcd}}}$. Tính $a+b+c+d$.
\shortans{$7$}
\loigiai{
Hình vuông $A_1$ có diện tích $S_1=1$.\\
Hình vuông $A_2$ là hình vuông có các đỉnh là trung điểm các cạnh của hình vuông $A_1$ do đó hình vuông $A_2$ có diện tích $S_2=\dfrac{1}{2}{S_1}=\dfrac{1}{2}$.\\
Tương tự, hình vuông $A_3$ có diện tích $S_3=\dfrac{1}{2}{S_2}=\dfrac{1}{2}\cdot\dfrac{1}{2}=\dfrac{1}{2^2}$.\\
Cứ tiếp tục như thế ta tính được diện tích hình vuông $A_{2024}$ là $S_{2024}=\dfrac{1}{2^{2023}}$.
}
\end{ex}

\begin{ex}%Câu 22%[1D2V2-7]
Một gia đình cần khoan một cái giếng để lấy nước. Họ thuê một đội khoan giếng nước đến để khoan giếng nước. Biết giá của một mét khoan đầu tiên là $75000$ đồng, kể từ mét khoan thứ hai giá của mỗi mét khoan tăng lên $6000$ đồng so với giá của mét khoan trước đó. Biết cần phải khoan sâu xuống $80$ m mới có nước. Số tiền phải trả (đơn vị triệu  đồng) là bao nhiêu?
\shortans{$25$}
\loigiai{
Số tiền phải trả để khoan giếng là một cấp số cộng có số hạng đầu $u_1=75000$; công sai $ d=6000$.\\
Số tiền phải trả để khoan $ 80m$ là $$S_{80}=\dfrac{80}{2}\left[2u_1+79d\right]=40\left[2.75000+79\cdot6000\right]=24\,960\,000 \, \text{đồng}\approx 25 \, \text{triệu đồng}.$$
}
\end{ex}
\Closesolutionfile{ans}
\indapan{6}{ans/ans-0-B15-KQ}